\section{Pure states, excision, and peaking projections}\label{sec:states}

The construction must preserve unique state extensions while changing selected
equivalence classes.  This section records the state-theoretic tools for doing
so.  Its main analytic fact is excision: near the support of a pure state,
compression of an element approaches its scalar state value.

All algebras and Hilbert spaces are complex.  Unless stated otherwise, algebras
are unital and inclusions preserve the unit.  A \emph{state} on an algebra $B$
is a positive linear functional of norm one, and it is \emph{pure} if it is an
extreme point of the state space.  Write $(\pi_h,H_h,\xi_h)$ for the GNS
representation of a state $h$.  Thus $\xi_h$ is a cyclic unit vector and
$h(a)=\langle\pi_h(a)\xi_h,\xi_h\rangle$; purity is equivalent to
irreducibility of $\pi_h$.  For pure states, write $h\sim k$ when $\pi_h$ and
$\pi_k$ are unitarily equivalent.  We use
$\Ad(u)(a)=uau^*$ for conjugation by a unitary.

A \emph{tracial algebra} is a unital $C^*$-algebra equipped with a faithful
tracial state $\tau$: one has $\tau(ab)=\tau(ba)$, and $\tau(a^*a)=0$ implies
$a=0$.  Traces and other bounded functionals are extended normally to the
enveloping von Neumann algebra $B^{**}$, with the same notation.  The normal
extension of a trace is tracial, by separate ultraweak continuity.
For a representation $\pi$ of $B$, its normal extension to $B^{**}$ is
denoted by $\pi^{**}$.

For an inclusion $B\subset C$, we shall use the property
\begin{equation}\tag{U}\label{eq:U}
 \text{every pure state of $B$ has exactly one state extension to $C$.}
\end{equation}
Here uniqueness is among all states.  States always extend across a unital
inclusion, by the positive form of the Hahn--Banach theorem.

\subsection{Unique extensions and peaking}

\begin{lemma}\label{na:unique-lift}
If a pure state $h$ of $B$ has a unique state extension $\widetilde h$ to
$C\supset B$, then $\widetilde h$ is pure.
\end{lemma}
\begin{proof}
In a nontrivial convex decomposition of $\widetilde h$, both restrictions to
$B$ equal $h$ by purity.  Uniqueness then makes both states equal to
$\widetilde h$.
\end{proof}

\begin{lemma}\label{na:unitary-states}
Equivalent pure states of a unital $C^*$-algebra $B$ are conjugate by a unitary
of $B$.  Consequently, if they each have a unique state extension to an
algebra $C\supset B$, their extensions are equivalent.
\end{lemma}
\begin{proof}
Realize the two states by unit vectors $\xi,\eta$ in a common irreducible
representation $\pi$.  On their finite-dimensional span choose a unitary
carrying $\xi$ to $\eta$ and a self-adjoint logarithm of that unitary.
Kadison transitivity, including its self-adjoint form
\cite[Theorem~1]{Kadison1957}, supplies a self-adjoint $a\in B$ whose action
on this span is the chosen logarithm.  The span reduces $\pi(a)$, so
$u=e^{ia}$ satisfies $\pi(u)\xi=\eta$.  If $h$ is represented by $\xi$
and $k$ by $\eta$, then $k=h\circ\Ad(u^*)$.
Conjugating their unique extensions by the same unitary proves the last
assertion.
\end{proof}

A decreasing sequence of projections $(P_n)_{n\ge0}$ in $B$ \emph{peaks to}
a state $h$ if $h$ is the only state satisfying $h(P_n)=1$ for every $n$.
Such a state is pure: both states in any convex decomposition satisfy the
same constraints.  The following elementary observation will let us pass to
smaller algebras without assuming \eqref{eq:U} for the inclusion.

\begin{lemma}[Peaking under restriction]\label{na:peaking-restriction}
Let $B\subset C$ be a unital inclusion and let $(P_n)$ be a decreasing
sequence of projections in $B$ that peaks to a state $H$ of $C$.
Then $(P_n)$ peaks to $h=H|_B$ in $B$, and $H$ is the unique state extension
of $h$ to $C$.
\end{lemma}
\begin{proof}
Every state of $B$ satisfying the projection constraints extends to a state
of $C$ with the same constraints, hence to $H$.  It must therefore equal
$h$.  Every extension of $h$ also satisfies the constraints and must equal
$H$.
\end{proof}

\subsection{Excision and inequivalent states}

For a state $h$, its \emph{support projection} $r_h\in B^{**}$ is the least
projection on which its normal extension takes value one.  We shall prove
that this projection is minimal when $h$ is pure, together with the norm
approximation that we need. The decreasing-net form of pure-state
excision used here is proved by Akemann, Anderson, and Pedersen
\cite[Proposition~2.2]{AkemannAndersonPedersen1986}. We record a proof
in the net form used below.

\begin{lemma}[Decreasing excision]\label{na:excision}
Let $h$ be a pure state of a unital $C^*$-algebra $B$.  There is a decreasing
net of positive contractions $x_\lambda\in B$ such that $h(x_\lambda)=1$,
$x_\lambda$ decreases strongly to the minimal projection $r_h$ in $B^{**}$,
and, for every $a\in B$,
\begin{equation}\label{eq:excision}
 \|x_\lambda(a-h(a)1)x_\lambda\|\longrightarrow0.
\end{equation}
Moreover, $r_h a r_h=h(a)r_h$ for every $a\in B$.
\end{lemma}
\begin{proof}
Put
\[
 L=\{a\in B:h(a^*a)=0\},\qquad N=L\cap L^*.
\]
Here $L$ is a closed left ideal and $N$ is a $C^*$-subalgebra.  If $a\in L$,
then $a^*a\in N$, since
$h((a^*a)^2)\le\|a\|^2h(a^*a)=0$.

First, $h$ is the only state annihilating $N$.  If another state $k$
annihilates $N$, then $L$ is contained in the kernel of
$a\mapsto\pi_k(a)\xi_k$.  The map
$Q_h:B\to H_h$, $Q_h(a)=\pi_h(a)\xi_h$, is onto by one-vector Kadison
transitivity.  This assertion concerns the whole Hilbert space, even when
it is nonseparable.  Since $\ker Q_h=L$, the open mapping theorem for
$B/L\to H_h$ shows that
\[
 T\bigl(\pi_h(a)\xi_h\bigr)=\pi_k(a)\xi_k
\]
defines a bounded operator $T:H_h\to H_k$.  It intertwines the
representations.  Thus $T^*T$ is scalar by irreducibility of $\pi_h$;
its value on $\xi_h$ shows that $T^*T=1$.  Taking vector coefficients gives
$k=h$.

We next construct the decreasing net.  Order $N_+$ by the usual operator
order, a directed order, and set
\[
 e_s=1-(1+s)^{-1},\qquad x_s=1-e_s \qquad(s\in N_+).
\]
Functional calculus gives $e_s\in N$, and operator monotonicity of
inversion makes $(e_s)$ increasing.  For $b\in N$ and
$s\ge m(bb^*+b^*b)$,
\[
 \|(1-e_s)b\|^2
 \le \frac1m\|(1+s)^{-1}s(1+s)^{-1}\|
 \le\frac1{4m};
\]
the same estimate holds for $\|b(1-e_s)\|^2$.  Thus $(e_s)$ is an
approximate identity of $N$, and $h(x_s)=1$.

Let $q$ be the strong limit of $(e_s)$, the support projection of $N$.
Every normal state supported on $1-q$ annihilates $N$, and hence is $h$;
conversely, $h(1-q)=1$.  The corner $(1-q)B^{**}(1-q)$ therefore has
exactly one normal state.  Normal states separate its self-adjoint
elements, so this corner is scalar.  It follows that $1-q=r_h$ is
minimal and $x_s\downarrow r_h$.  The identity $r_har_h=h(a)r_h$ follows
by applying $h$ to this scalar corner.

Finally, let $\rho_s$ be arbitrary states of $B$.  Every weak-star cluster
point of the positive functionals
$a\mapsto\rho_s(x_sax_s)$ annihilates $N$, by the approximate-identity
estimates.  A nonzero cluster point, after normalization, is therefore
$h$.  For self-adjoint $a$, a failure of \eqref{eq:excision} would give a
subnet and norming states $\rho_s$ with a cluster point nonzero on
$a-h(a)1$, a contradiction.  Real and imaginary parts give the conclusion
for arbitrary $a$.
\end{proof}

\begin{lemma}[Projection excision]\label{na:projection-excision}
If decreasing projections $(P_n)$ in $B$ peak to $h$, then
$P_n\downarrow r_h$ strongly and
\[
 \|P_n(a-h(a)1)P_n\|\longrightarrow0\qquad(a\in B).
\]
In particular, if $h(a)=0$, then $\|P_n a P_t\|\to0$ on the product
index set.
\end{lemma}
\begin{proof}
The infimum of the $P_n$ in $B^{**}$ supports exactly one normal state,
namely $h$.  As in the preceding proof, it is the minimal projection
$r_h$.  For arbitrary states $\rho_n$, every cluster point $\psi$ of
$\rho_n(P_n\,\cdot\,P_n)$ vanishes on $1-P_m$ for each fixed $m$.
If $\psi\ne0$, its normalization satisfies every peaking constraint and
therefore equals $h$.  The same norming-state argument proves excision.
For the last assertion use nesting:
\[
 \|P_n a P_t\|\le \|P_{\min(n,t)}aP_{\min(n,t)}\|.
\qedhere
\]
\end{proof}

Excision also separates inequivalent states in norm.  We use the following
form of the cross-excision estimate of Akemann and Weaver
\cite[Lemma~1]{AkemannWeaver2004}.

\begin{lemma}[Cross-excision]\label{na:cross-excision}
Let $h\not\sim k$ be pure states of $B$, and let $(x_\lambda)$ and
$(y_\nu)$ be decreasing nets as in Lemma~\ref{na:excision}, for $h$ and
$k$ respectively.  Then, for every $a\in B$,
\begin{equation}\label{eq:cross-excision}
 \|x_\lambda a y_\nu\|\longrightarrow0
\end{equation}
on the product directed set.  Either net may be a peaking sequence of
projections.
\end{lemma}
\begin{proof}
The normal extension of $\pi_h$ annihilates $r_k$.  Otherwise a unit vector
in the range of $\pi_h^{**}(r_k)$ would represent $k$, by
$r_kar_k=k(a)r_k$.  Every nonzero vector in an irreducible representation
is cyclic, so this would imply $h\sim k$.

Consequently $h(ay_{\nu_0}a^*)\to0$ as $\nu_0$ increases.  For
$\nu\ge\nu_0$, the inequalities $y_\nu^2\le y_\nu\le y_{\nu_0}$ give
\begin{align*}
 \|x_\lambda a y_\nu\|^2
 &\le \|x_\lambda a y_{\nu_0}a^*x_\lambda\|\\
 &\le h(ay_{\nu_0}a^*)
   +\|x_\lambda(ay_{\nu_0}a^*-h(ay_{\nu_0}a^*)1)x_\lambda\|.
\end{align*}
Choose $\nu_0$ first, then use excision to choose $\lambda_0$.
This proves smallness on a rectangular tail and hence the product-net
limit.  Lemma~\ref{na:projection-excision} gives the stated projection
case.
\end{proof}

\begin{lemma}[Trace-null supports]\label{na:trace-null}
If $B$ has no finite-dimensional irreducible representations and $\tau$ is
a tracial state on $B$, then $\tau(r_h)=0$ for every pure state $h$.
\end{lemma}
\begin{proof}
The space $H_h$ is infinite-dimensional.  Choose orthonormal unit vectors
$\eta_j$ and unitaries $u_j\in B$ with $\pi_h(u_j)\xi_h=\eta_j$, as in
Lemma~\ref{na:unitary-states}.  The projections $u_jr_hu_j^*$ are
pairwise orthogonal, since
\[
 r_hu_i^*u_jr_h=h(u_i^*u_j)r_h=0\qquad(i\ne j).
\]
They all have trace $\tau(r_h)$, so $m\tau(r_h)\le1$ for every positive
integer $m$.
\end{proof}

\subsection{Expectations from peaking projections}

We now use projection excision to construct the conditional expectations
needed for the tracial extensions.  A \emph{conditional expectation}
$E:C\to B$ is a unital completely positive map fixing $B$ and satisfying
$E(b_1ab_2)=b_1E(a)b_2$ for $b_1,b_2\in B$.  It is \emph{faithful} if
$E(a^*a)=0$ implies $a=0$.

\begin{lemma}\label{na:trace-expectation}
For a unital subalgebra $B$ of a faithfully tracial algebra $(C,\tau)$,
there is at most one trace-preserving conditional expectation $C\to B$.
If one exists, it is faithful.
\end{lemma}
\begin{proof}
For such an expectation, bimodularity gives
$\tau(b^*E(a))=\tau(b^*a)$ for every $b\in B$.  Two expectations therefore
have difference $d\in B$ satisfying $\tau(d^*d)=0$, so they agree.
Faithfulness follows from $\tau(E(a^*a))=\tau(a^*a)$.
\end{proof}

\begin{lemma}\label{na:peaking-expectation}
Let $(B,\tau)$ be faithfully tracial.  Suppose decreasing projections
$(p_n)_{n\ge0}$ peak to $h$, with $p_0=1$ and trace values strictly
decreasing to zero.  Then
$D=C^*(p_n:n\ge0)$ is isomorphic to
$C(\mathbb N_0\cup\{\infty\})$, and there is a trace-preserving conditional
expectation $E:B\to D$.  Writing $d_n=p_n-p_{n+1}$ for the finite-point
atoms, it is given by
\begin{equation}\label{eq:peaking-expectation}
 E(a)(n)=\frac{\tau(d_nad_n)}{\tau(d_n)},\qquad E(a)(\infty)=h(a).
\end{equation}
\end{lemma}
\begin{proof}
The $d_n$ are nonzero orthogonal projections.  Their closed linear span is
isomorphic to $c_0(\mathbb N_0)$, and adjoining the unit gives the claimed
description of $D$.  Indeed $p_n=1-\sum_{j<n}d_j$ and $\|p_n\|=1$,
since $\tau(p_n)>0$.  Multiplying by $p_n$ shows that $1$ has distance
at least one from every linear combination of $d_0,\ldots,d_{n-1}$.

The values in \eqref{eq:peaking-expectation} define a continuous function:
by $d_n\le p_n$ and Lemma~\ref{na:projection-excision},
\[
 \left|\frac{\tau(d_nad_n)}{\tau(d_n)}-h(a)\right|
 \le \|p_n(a-h(a)1)p_n\|\longrightarrow0.
\]
Every coordinate of $E$ is a state, so $E$ is unital and completely
positive: positivity of each matrix amplification is tested pointwise
in the commutative range.  The map fixes $D$.  To check bimodularity,
an element of $D$ acts scalarly on each $d_n$; at infinity, $h|_D$ is the
evaluation character.  Cauchy--Schwarz for $h$ then gives
$h(ba)=h(b)h(a)$ and $h(ab)=h(a)h(b)$ for $b\in D$.
These coordinate identities prove bimodularity.

Finally, $\sum_n\tau(d_n)=1$ and the trace on $D$ has no mass at infinity.
Summing the finite coordinates gives $\tau(E(a))=\tau(a)$, since the
remaining trace satisfies
$|\tau(p_Na)|\le\|a\|\tau(p_N)\to0$.
\end{proof}

Unique state extensions and a faithful expectation will also preserve
simplicity, without a separate analysis of ideals in each extension.

\begin{lemma}\label{na:simple-extension}
Suppose $B\subset C$ satisfies \eqref{eq:U}, $B$ is simple, and there is a
faithful conditional expectation $E:C\to B$.  Then $C$ is simple.
\end{lemma}
\begin{proof}
Let $J$ be a proper closed two-sided ideal of $C$.  The quotient map
$q:C\to C/J$ is injective on $B$, since $B$ is simple and $q(1)\ne0$.
Every pure state $h$ of $B$ therefore extends through this quotient to a
state of $C$ annihilating $J$.  By \eqref{eq:U}, this extension is $h\circ E$.
Thus $h(E(a))=0$ for every positive $a\in J$ and every pure $h$ on $B$.
Pure states separate positive elements, so $E(a)=0$; faithfulness gives
$a=0$.  Hence $J=0$.
\end{proof}
